\section*{第 \ref{ch:g11:trigo} 章——三角函数：单位圆}

\begin{solution}{exo:g11:trigo:1}
化为\omterm{def:g11:trigo:radian}{弧度}：$30^\circ = \frac{\pi}{6}$，$45^\circ = \frac{\pi}{4}$，
$120^\circ = \frac{2\pi}{3}$, $270^\circ = \frac{3\pi}{2}$.

化为角度：$\frac{\pi}{5} = 36^\circ$，$\frac{5\pi}{6} = 150^\circ$，
$\frac{7\pi}{4} = 315^\circ$, $\frac{2\pi}{9} = 40^\circ$.
\end{solution}

\begin{solution}{exo:g11:trigo:2}
$\frac{2\pi}{3}$ 与 $-\frac{\pi}{4}$ 已在
$\intoc{-\pi}{\pi}$ 内：分别位于第二象限和第四象限。

$\frac{17\pi}{6} - 2\pi = \frac{5\pi}{6}$：位于第二象限，靠近
负 $x$ 轴。

$-\frac{7\pi}{2} + 4\pi = \frac{\pi}{2}$：位于单位圆顶端，
$(0, 1)$。
\end{solution}

\begin{solution}{exo:g11:trigo:3}
$\cos\frac{3\pi}{4} = \cos\left(\pi - \frac\pi4\right)
= -\cos\frac\pi4 = -\frac{\sqrt2}{2}$.

$\sin\frac{5\pi}{6} = \sin\left(\pi - \frac\pi6\right)
= \sin\frac\pi6 = \frac12$.

$\cos\left(-\frac\pi3\right) = \cos\frac\pi3 = \frac12$.

$\sin\frac{4\pi}{3} = \sin\left(\pi + \frac\pi3\right)
= -\sin\frac\pi3 = -\frac{\sqrt3}{2}$.

$\cos\frac{11\pi}{6} = \cos\left(-\frac\pi6 + 2\pi\right)
= \cos\frac\pi6 = \frac{\sqrt3}{2}$.
\end{solution}

\begin{solution}{exo:g11:trigo:4}
$\sin^2 t = 1 - \cos^2 t = 1 - \frac{25}{169} = \frac{144}{169}$，于是
$\sin t = \pm\frac{12}{13}$。在 $\intoo{-\frac\pi2}{0}$ 内（第四象限）
\omterm{def:g10:coordgeom:system}{纵坐标}为负，故 $\sin t = -\frac{12}{13}$。
\end{solution}

\begin{solution}{exo:g11:trigo:5}
由 \cref{prop:g11:trigo:associated} 得：
\[
A = (-\cos x) + \cos x + \cos x - \cos x = 0 ,
\]
\[
B = \sin x + (-\sin x) + (-\sin x) = -\sin x .
\]
\end{solution}

\begin{solution}{exo:g11:trigo:6}
$\cos t = \frac{\sqrt3}{2}$：一个解为 $\frac\pi6$，故
$t = \pm\frac\pi6$（$2k\pi$ 的其他平移均超出该\omterm{def:g10:numbers:interval}{区间}）。

$\sin t = \frac12$：$t = \frac\pi6$ 或 $t = \pi - \frac\pi6 =
\frac{5\pi}{6}$.

$\cos t = -1$：\omterm{def:g10:coordgeom:system}{横坐标}为 $-1$ 的点只有 $M(\pi)$，故 $t = \pi$。
\end{solution}

\begin{solution}{exo:g11:trigo:7}
$\sin t = -\frac{\sqrt3}{2}$：由 $a = -\frac\pi3$，两族解分别为
$-\frac\pi3 + 2k\pi$ 与 $\pi + \frac\pi3 + 2k\pi$；在
$\intco{0}{2\pi}$ 内：$t = \frac{4\pi}{3}$ 与 $t = \frac{5\pi}{3}$。

$\cos t = 0$：$t = \frac\pi2$ 与 $t = \frac{3\pi}{2}$。

$2\sin t + 1 = 0$ 即 $\sin t = -\frac12$：由 $a = -\frac\pi6$，在
$\intco{0}{2\pi}$ 内：$t = \frac{7\pi}{6}$ 与 $t = \frac{11\pi}{6}$。
\end{solution}

\begin{solution}{exo:g11:trigo:8}
由 \cref{thm:g11:trigo:equations} 可知，$\cos 2t = \cos t$ 等价于
$2t = t + 2k\pi$ 或 $2t = -t + 2k\pi$，即 $t = 2k\pi$ 或
$t = \frac{2k\pi}{3}$（$k \in \Z$）。第一族包含于
第二族中（取 $k$ 为 $3$ 的倍数），因此解集为
\[
S = \left\{\frac{2k\pi}{3} : k \in \Z\right\}
\]
——即单位圆上的三个点 $M(0)$，
$M\!\left(\frac{2\pi}{3}\right)$, $M\!\left(\frac{4\pi}{3}\right)$.
\end{solution}

\begin{solution}{exo:g11:trigo:9}
\omterm{def:g10:algebra:expand}{展开}两个平方并利用 $\cos^2 x + \sin^2 x = 1$：
\[
(\cos^2 x + 2\sin x\cos x + \sin^2 x)
+ (\cos^2 x - 2\sin x\cos x + \sin^2 x)
= 1 + 1 = 2 .
\]
\end{solution}

\begin{solution}{exo:g11:trigo:10}
无滑动滚动是指弧长等于
行进的距离，即 $r\theta = 150$ cm，由此得
$\theta = \frac{150}{30} = 5$ \omterm{def:g11:trigo:radian}{弧度} $= 5 \times \frac{180}{\pi} \approx
286.5^\circ$。车轮转一整圈，车子前进的圆周长
$2\pi \times 30 = 60\pi \approx 188.5$ cm，约 $1.88$ m。
\end{solution}

\begin{solution}{exo:g11:trigo:11}
在单位圆上，\omterm{def:g10:coordgeom:system}{横坐标}恰为 $\frac12$ 的点为
$M\!\left(\pm\frac\pi3\right)$。\omterm{def:g10:coordgeom:system}{横坐标}\emph{不超过}
$\frac12$ 的点构成竖直直线 $x = \frac12$ 左侧的弧，
当 $t$ 从 $\frac\pi3$ 增至 $\pi$，或从 $-\pi$ 增至
$-\frac\pi3$ 时所经过的弧。因此，在 $\intoc{-\pi}{\pi}$ 上：
\[
S = \intoc{-\pi}{-\frac{\pi}{3}} \cup \intcc{\frac{\pi}{3}}{\pi}.
\]
（因不等式为非严格不等式，端点 $\pm\frac\pi3$ 均包含在内。）
\end{solution}

\begin{solution}{pb:g11:trigo:1}
\textbf{1.} $30^\circ = \frac{\pi}{6}$；$135^\circ =
\frac{3\pi}{4}$；$300^\circ = \frac{5\pi}{3}$。又
$\frac{3\pi}{4} = 135^\circ$；$\frac{5\pi}{6} = 150^\circ$。

\textbf{2.} 弧长 $r\,t = 3 \times 2 = 6$ 米。用\omterm{def:g11:trigo:radian}{弧度}时，弧长
为简单的乘积——若用角度，每个公式里都要带上
一个 $\frac{\pi}{180}$ 因子。

\textbf{3.} 转过的\omterm{def:g11:trigo:radian}{弧度}数 $= \frac{\text{距离}}
{\text{半径}} = \frac{1000}{0.30} \approx 3\,333$ \omterm{def:g11:trigo:radian}{弧度}；
除以 $2\pi$：约 $530$ 整圈。

\textbf{4.} $\cos\frac{2\pi}{3} = -\frac12$；
$\sin\left(-\frac\pi4\right) = -\frac{\sqrt2}{2}$；
$\frac{19\pi}{6} - 2\pi = \frac{7\pi}{6}$，于是
$\cos\frac{19\pi}{6} = \cos\frac{7\pi}{6} =
-\frac{\sqrt3}{2}$。

\textbf{5.} $\cos^2 t = 1 - \frac{9}{25} = \frac{16}{25}$；在
第二象限中\omterm{def:g11:trigo:cossin}{余弦}为负：
$\cos t = -\frac45$，且
$\tan t = \frac{3/5}{-4/5} = -\frac34$。

\textbf{6.} 经过 $t$ 分钟，摩天轮转过的角为
$\frac{t}{30} \times 2\pi = \frac{\pi t}{15}$。当 $t = 0$ 时：
$h = 65 - 60 = 5$ 米——底部的登车平台
（轮心高度减去半径）。负号使你\emph{低于}
轮心，此时转过的角为 $0$。

\textbf{7.} $h(7.5) = 65 - 60\cos\frac{\pi}{2} = 65$ 米（轮心
高度）；$h(15) = 65 + 60 = 125$ 米（最高点）；
$h(20) = 65 - 60\cos\frac{4\pi}{3} = 65 + 30 = 95$ 米。

\textbf{8.} 由 $65 - 60\cos\frac{\pi t}{15} = 95$ 得
$\cos\frac{\pi t}{15} = -\frac12$：第一个解为
$\frac{\pi t}{15} = \frac{2\pi}{3}$，即 $t = 10$ 分钟。

\textbf{9.} $\cos\frac{\pi t}{15} \leq -\frac12$ 在
$\frac{\pi t}{15} \in \intcc{\frac{2\pi}{3}}{\frac{4\pi}{3}}$ 时成立，
即 $t \in \intcc{10}{20}$：乘坐时间中有十分钟高于
$95$ 米，关于 $t = 15$ 时的最高点对称。

\textbf{10.} 第一分钟：
\[
h(1) - h(0) = 60\left(1 - \cos\frac{\pi}{15}\right)
\approx 1.3 \text{ 米}.
\]
第 $7$ 与 $8$ 分钟之间：
\[
60\left(\cos\frac{7\pi}{15} - \cos\frac{8\pi}{15}\right)
\approx 12.5 \text{ 米}
\]
——几乎是前者的十倍。高度变化最快出现在经过轮心高度时
而在接近最低点与最高点处变化趋缓：\omterm{def:g11:trigo:cossin}{正弦}曲线在
中部陡峭，在两端平缓。

\textbf{11.} 半径 $= \frac{40 - 4}{2} = 18$ 米，轮心位于
$4 + 18 = 22$ 米，周期为 $12$ 分钟：
$h(t) = 22 - 18\cos\left(\frac{\pi t}{6}\right)$。

\textbf{12.} \omterm{def:g10:functions:extrema}{最大值}为 $7 + 3 = 10$ 米，当
$\sin\frac{\pi t}{6} = 1$ 时，首次出现在 $t = 3$（凌晨 3 点）；
\omterm{def:g10:functions:extrema}{最小值} $4$ 米，首次出现在 $t = 9$；周期为
$\frac{2\pi}{\pi/6} = 12$ 小时。

\textbf{13.} $\sin\frac{\pi t}{6} \geq \frac12$ 在
$\frac{\pi t}{6} \in \intcc{\frac\pi6}{\frac{5\pi}{6}}$ 时成立：
即 $t \in \intcc{1}{5}$——一段四小时的时间，从凌晨 1 点到
凌晨 5 点，一个周期后再次开放，$t \in \intcc{13}{17}$
（下午 1 点至下午 5 点）。

\textbf{14.} 关于纵轴的镜面反射将角为
$t$ 的点映到角为 $\pi - t$ 的点，并保持
\omterm{def:g10:coordgeom:system}{纵坐标}不变：$\sin(\pi - t) = \sin t$。绕
圆心的半圈旋转将 $t$ 变为 $\pi + t$，两坐标均取相反数：
$\cos(\pi + t) = -\cos t$。而关于对角线的镜面反射
交换\omterm{def:g10:coordgeom:system}{横坐标}与\omterm{def:g10:coordgeom:system}{纵坐标}：
$\cos\left(\frac\pi2 - t\right) = \sin t$——正是
余角关系中“co”的含义。

\textbf{15.} $\sin(-t) = -\sin t$：\omterm{def:g11:func:parity}{奇函数}（\omterm{def:g10:functions:graph}{图像}关于原点
中心对称）；$\cos(-t) = \cos t$：\omterm{def:g11:func:parity}{偶函数}（镜面对称）。而
$\sin t + t^3$ 是两个\omterm{def:g11:func:parity}{奇函数}之和，故为\omterm{def:g11:func:parity}{奇函数}。

\textbf{16.} \omterm{def:g10:algebra:expand}{因式分解} $(2\sin t + 1)(\sin t - 1) = 0$：
$\sin t = 1$ 得 $t = \frac\pi2$；
$\sin t = -\frac12$ 得 $t = \frac{7\pi}{6}$ 与
$t = \frac{11\pi}{6}$。单位圆上共有三个解。

\textbf{17.} $1$ \omterm{def:g11:trigo:radian}{弧度} $= \frac{180}{\pi} \approx 57.3^\circ$。
因此 $\sin(1) \approx 0.841$，而 $\sin(1^\circ) \approx
0.0175$：相差近 $50$ 倍。教训：先检查模式指示灯，
再相信任何三角\omterm{def:g11:func:function}{函数}的显示结果。

\textbf{18.} 对角线上\omterm{def:g10:coordgeom:system}{横坐标}与\omterm{def:g10:coordgeom:system}{纵坐标}相等：
$t = \frac\pi4$ 与 $t = \frac{5\pi}{4}$。

\textbf{19.} $\sin(0.1) = 0.09983\ldots$ 与 $0.1$ 相比：
相对误差约为 $0.17\,\%$。对于小角，弦
紧贴弧线——而十二年级的极限
$\frac{\sin t}{t} \to 1$ 正是把这一观察
升华为定理。

\textbf{20.} \omterm{def:g11:trigo:radian}{弧度}制下角 $=$ 单位圆上的弧长，
因此角的运算与长度相同。\omterm{def:g11:trigo:cossin}{余弦}与\omterm{def:g11:trigo:cossin}{正弦}是
时针端点的坐标——地面上的影子，墙上的
高度。圆的对称性：一整套
恒等式，可由镜面与半圈旋转直接读出。\omterm{def:g10:algebra:equation}{方程}
$\cos t = c$：钟面上的时刻表，借助单位圆
一次看清所有解。装束：摩天轮穿着
$65 - 60\cos$，潮汐穿着 $7 + 3\sin$——同一只时钟，
同一种数学。
\end{solution}
